\documentclass[10pt,a4paper]{amsart}
\usepackage[margin=19mm]{geometry}
\usepackage{amsmath,amssymb,mathtools}
\usepackage[scr=rsfs,cal=euler]{mathalfa}
\usepackage{xfrac}
\usepackage[unicode,hidelinks,bookmarks=false]{hyperref}
\newcommand{\ie}{i.e.\ }
\newcommand{\cf}{cf.\ }
\newcommand{\resp}{respectively\ }
\newcommand{\Subsection}{\subsection}
\providecommand{\nobreakdash}{\nobreak\hbox{-}}
\setlength{\emergencystretch}{2em}
\multlinegap0pt
\usepackage{amssymb}
\usepackage{xcolor}
\usepackage[shortlabels]{enumitem}
\usepackage{mathtools}
\newtheorem{lemma}{Lemma}
\newtheorem{theorem}{Theorem}
\newcommand{\SSD}{\operatorname{SSD}}
\newcommand{\Z}{\mathbb Z}
\title{Semisimple derivations, rational slices, and kernels over affine domains}
\author{Luis Cid}
\date{Source: arXiv:2603.19587v1 (CC BY 4.0)}
\begin{document}
\maketitle

\noindent\emph{Source and licence.} Semisimple derivations, rational slices, and kernels over affine domains. Version arXiv:2603.19587v1 (CC BY 4.0), distributed by arXiv under the Creative Commons Attribution 4.0 International licence, http://creativecommons.org/licenses/by/4.0/, as recorded in the arXiv licence metadata for this exact version and shown on the abstract page https://arxiv.org/abs/2603.19587v1. Full licence notice: \texttt{licenses/arxiv-2603.19587-CC-BY-4.0.txt}.\par\smallskip

\noindent\textit{Editorial scope.} Counted unit: the article's theorem thm:kernel-fraction-field (source0.tex lines 552--568). The article's complete original proof of it is delivered with it (source0.tex lines 570--656, one \texttt{proof} environment of 87 lines). No mathematics is written, repaired or re-derived: the statement and the proof are the article's own lines, in the article's own notation, and no retained line is substituted or re-flowed.  Retained uncounted as the source-local context the proof needs: lines 403--416.  Nothing was omitted from any retained span, so the package carries no omission record.  No retained line contains a citation, so the wrapper carries no bibliography and no bibliography entry.  No retained line contains a figure environment, an image-inclusion command or drawing code, so no image is delivered and the package declares no asset dependency.  Editorial formatting only, in addition: an AMS \texttt{amsart} preamble that restates the article's own declarations and macros used by the retained lines, namely the environments \texttt{lemma}, \texttt{theorem} on separate counters, together with the standard packages those lines need.  The retained environments print in this document as Lemma 1, Theorem 1 in order of appearance, whereas the article prints the same items with the numbering of its own full document.  The complete native source of the article as distributed in the arXiv e-print for this exact version is bundled unchanged as \texttt{sources/196786/source0.tex}.
\begin{lemma}\label{lem:kernel-weighted-element}
Let $D\in \SSD(B)$ and let $s\in K$ satisfy
\[
D(s)=s.
\]
If $b\in K$ is a semi-invariant of weight $\lambda$, that is,
\[
D(b)=\lambda b,
\]
with $\lambda\in \Z$, then
\[
bs^{-\lambda}\in \ker(D|_K).
\]
\end{lemma}

\begin{theorem}\label{thm:kernel-fraction-field}
Assume that $D\in \SSD^{*}(B)$, and let $s\in K$ be such that $D(s)=s$.
Suppose that $K$ is generated over $k$ by semi-invariants
\[
b_1,\dots,b_n\in K,
\qquad
D(b_i)=\lambda_i b_i,
\qquad
\lambda_i\in \Z.
\]
Then
\[
\ker(D)=k\left(b_1s^{-\lambda_1},\dots,b_ns^{-\lambda_n}\right),
\qquad
K=\ker(D)(s).
\]
\end{theorem}

\begin{proof}
Set
\[
u_i:=b_i s^{-\lambda_i}\qquad (i=1,\dots,n).
\]
By Lemma~\ref{lem:kernel-weighted-element}, each $u_i$ belongs to $\ker(D)$. Hence
\[
L:=k(u_1,\dots,u_n)\subseteq \ker(D).
\]

On the other hand, for each $i$ we have
\[
b_i=u_i s^{\lambda_i},
\]
so
\[
K=k(b_1,\dots,b_n)\subseteq L(s)\subseteq K.
\]
Therefore
\[
K=L(s).
\]

It remains to prove that
\[
\ker(D)=L.
\]

We claim that $s$ is transcendental over $L$. Indeed, assume that $s$ is algebraic over $L$.
Choose a nonzero polynomial
\[
P(T)=a_0+a_1T+\cdots+a_mT^m\in L[T]
\]
of minimal degree $m$ such that
\[
P(s)=0.
\]
Since $D$ vanishes on $L$ and satisfies $D(s)=s$, applying $D$ to $P(s)=0$ gives
\[
0=D(P(s))=\sum_{i=1}^m i\,a_i s^i.
\]
Because $s\neq 0$ in the field $K$, we can factor out $s$:
\[
0=\sum_{i=1}^m i\,a_i s^{i-1} = a_1 + 2a_2 s + \cdots + m a_m s^{m-1}.
\]
This is a polynomial relation over $L$ of degree $m-1$ satisfied by $s$, and its leading
coefficient $m a_m$ is nonzero since $\operatorname{char}(k)=0$ and $a_m\neq 0$.
This contradicts the minimality of $m$. Hence $s$ is transcendental over $L$.

Therefore $K=L(s)$ is a rational function field in one variable over $L$, and $D$ acts as the Euler derivation with respect to the parameter $s$, namely
\[
D=s\frac{d}{ds}.
\]
Let $f\in \ker(D)$. Since $K=L(s)$, we may write
\[
f=\frac{P(s)}{Q(s)}
\]
with $P,Q\in L[s]$ and $Q\neq 0$.

Thus
\[
0=D(f)=s\frac{df}{ds}.
\]
Since $K$ is a field and $s\neq 0$, this implies
\[
\frac{df}{ds}=0.
\]
In characteristic zero, the only rational functions in $L(s)$ with zero derivative are the elements of $L$. Hence $f\in L$. Therefore
\[
\ker(D)=L
=
k(u_1,\dots,u_n)
=
k\left(b_1s^{-\lambda_1},\dots,b_ns^{-\lambda_n}\right).
\]

The equality
\[
K=\ker(D)(s)
\]
now follows from
\[
K=L(s)
\qquad\text{and}\qquad
L=\ker(D).
\]
\end{proof}

\end{document}
